\scanpage{24}
Soit $\mathcal E$ la catégorie des ensembles finis, avec comme
morphismes les applications \uline{injectives}. On
considère des foncteurs
\[
\mathcal E\xrightarrow{F}(\mathrm{Ens})
\]
tels que $\forall E\in\operatorname{Ob}\mathcal E$, $\operatorname{card}(F(E))$ est fini, i.e. des
foncteurs
\[
\mathcal E\longrightarrow(\mathrm{Ens}_f)
\]
Ils forment une catégorie $\uline{\operatorname{Hom}}(\mathcal E,(\mathrm{Ens}_f))$,
qui est un para-topos. Parmi ses objets, on
en distingue deux sortes particulières

a) Les foncteurs \uline{représentables} $\mathcal E\mathbin{\text{---}}(\mathrm{Ens})$.
de la forme
\[
h_I:E\longmapsto\operatorname{Hom}_{\mathcal E}(I,E)=\operatorname{Mon}(I,E)
\]
où $I$ est un ensemble fini. Le foncteur
\[
I\longrightarrow h_I\qquad\mathcal E^o\longrightarrow\uline{\operatorname{Hom}}(\mathcal E,(\mathrm{Ens}_f))
\]
est pleinement fidèle, et identifie $\mathcal E^o$ à
une sous-catégorie pleine de $\uline{\operatorname{Hom}}(\mathcal E,(\mathrm{Ens}_f))$

b) Les foncteurs de la forme
\[
H_I:E\longrightarrow\operatorname{Hom}_{(\mathrm{Ens})}(I,E)=E^I=H(I,E)
\]
(restrictions à $\mathcal E\subset(\mathrm{Ens}_f)$ des foncteurs
correspondants sur $(\mathrm{Ens}_f)$).

\uline{Prop} Les $h_I\in\operatorname{Ob}\overbrace{\uline{\operatorname{Hom}}(\mathcal E,(\mathrm{Ens}_f))}^{\mathcal H}$ sont des
objets connexes (non vides) de $\mathcal H$. À isomorphisme
\mspage{23}
